\chapter{Alessio Figalli}\label{ch2018alessiofigalli:alessio-figalli-fields-medal-2018-a-comprehensive-report-on-his-work-and-impact}

\begin{center}\rule{0.5\linewidth}{0.5pt}\end{center}

\textbf{Prepared:} July 2026

\textbf{Subject:} Alessio Figalli (b.\ April 2, 1984, Rome, Italy)

\textbf{Position:} Chaired Professor of Mathematics, ETH Zurich; Director, Forschungsinstitut f\"ur Mathematik (FIM)

\textbf{Doctoral Advisors:} Luigi Ambrosio (Scuola Normale Superiore di Pisa), C\'edric Villani (ENS Lyon)

\textbf{Fields Medal Citation:} \emph{``For contributions to the theory of optimal transport, and its applications to partial differential equations, metric geometry, and probability.''}

\begin{center}\rule{0.5\linewidth}{0.5pt}\end{center}

\begin{prereqs}
Optimal transport (Monge problem, Wasserstein distance, Monge--Amp\`ere equation), PDE theory (elliptic equations, free boundary problems), basic measure theory and Riemannian geometry.
\end{prereqs}

\begin{objectives}
Explain the regularity theory \cite{figalli2016} for optimal transport maps and the Monge--Amp\`ere equation; describe quantitative stability results for isoperimetric and functional inequalities; understand applications to free boundary problems and the Keller--Segel equation.
\end{objectives}

\section{Problem History}

The optimal transport problem, first formulated by Gaspard Monge in 1781, asks: given two probability distributions $\mu$ and $\nu$ on measurable sets $\Omega_1, \Omega_2 \subset \mathbb{R}^n$, find a map $T: \Omega_1 \to \Omega_2$ transporting $\mu$ to $\nu$ while minimizing $\int_{\Omega_1} c(x, T(x))\,d\mu(x)$ for some cost function $c$. For the quadratic cost $c(x,y) = |x-y|^2$, Brenier (1991) showed that the optimal map is the gradient of a convex function $T = \nabla\phi$. This was the first major regularity result in optimal transport theory, connecting the Monge problem to the theory of convex functions and fully nonlinear PDEs.

The function $\phi$ satisfies the \textbf{Monge--Amp\`ere equation}: $\det(D^2\phi(x)) = \frac{d\mu}{d\nu}(\nabla\phi(x))$. Regularity theory---understanding when $\phi$ is smooth---was pioneered by Caffarelli in the 1990s for interior points with convex domains and bounded, smooth densities. Boundary regularity, non-convex domains, and rough target densities remained challenging. The Monge--Amp\`ere equation is a fully nonlinear PDE that arises naturally in differential geometry (the Minkowski problem), meteorology (the semi-geostrophic equations), and economics (the problem of optimal allocation of resources).

The \textbf{partial regularity problem}---characterizing the singular set where $\nabla\phi$ fails to be continuous---was a major open question. The singular set could a priori be very large, and understanding its structure required new techniques beyond those available from classical PDE theory. This problem is closely related to the question of uniqueness of optimal transport maps and to the structure of the space of probability measures.

The \textbf{isoperimetric inequality} has sharp equality cases (the sphere), but \textbf{stability}---bounding the distance to the extremal in terms of the deficit---was largely unknown for general norms and in non-Euclidean settings. Stability results are essential for numerical analysis, geometric flows, and applications to PDEs.

Free boundary problems, where the domain of a PDE is unknown, arise from optimal transport (partial transport), the obstacle problem, and the Stefan problem. Regularity of the free boundary is a fundamental challenge in modern PDE theory. The connection between optimal transport and free boundaries was discovered in the study of the partial transport problem.

The \textbf{Keller--Segel equation} models chemotaxis (the aggregation of cells in response to chemical signals). For critical initial mass, the question of whether solutions blow up in finite time was a major open problem. The connection to the logarithmic Sobolev inequality and Wasserstein gradient flows provided a new approach. The Keller--Segel equation is a parabolic-elliptic PDE system that arises in mathematical biology and has been studied extensively by mathematicians and biologists alike.

The \textbf{semi-geostrophic equations} in meteorology describe the motion of air masses under the influence of the Coriolis force. The equations have a variational structure that is naturally described by optimal transport theory. Figalli's work on the regularity of solutions to these equations has applications to weather prediction and climate modeling, demonstrating the deep connections between pure mathematics and applied science.

The \textbf{Wasserstein distance} $W_2(\mu, \nu)$ between two probability measures is defined as the infimum of $\int |x - T(x)|^2\,d\mu(x)$ over all maps $T$ transporting $\mu$ to $\nu$. This distance metrizes the weak convergence of probability measures and provides a natural geometry for the space of probability measures. The regularity of optimal transport maps is the key to understanding this geometry. The Wasserstein distance has become a fundamental tool in probability theory, statistics, and machine learning, with applications ranging from image processing to natural language processing.

The \textbf{Benamou--Brenier formula} expresses the Wasserstein distance as an infimum over time-dependent vector fields: $W_2(\mu, \nu)^2 = \inf \int_0^1 \int |v_t|^2\,d\mu_t\,dt$, where the infimum is over all pairs $(v_t, \mu_t)$ satisfying the continuity equation $\partial_t \mu_t + \nabla \cdot (v_t \mu_t) = 0$. This formula connects optimal transport to the theory of fluid dynamics and provides a variational characterization of the Wasserstein distance that is useful for numerical computations.

The \textbf{Monge--Amp\`ere equation} is a fully nonlinear PDE that arises naturally in differential geometry (the Minkowski problem), meteorology (the semi-geostrophic equations), and economics (the problem of optimal allocation of resources). Regularity theory for this equation is one of the central problems in nonlinear PDE theory, and Figalli's work has made significant progress toward a complete understanding of the regularity of solutions.

\section{How Was It Solved}

\textbf{Partial regularity for Monge--Amp\`ere (2010--2012).} With De Philippis, Figalli proved that for a convex solution of $\det D^2 u = f$ with $0 < \lambda \le f \le \Lambda$, the singular set has Hausdorff dimension at most $n-1$. The optimal transport map $T = \nabla u$ is therefore continuous outside a set of codimension at least 1. The proof uses a careful analysis of the geometry of level sets and the integrability of second derivatives.

The key innovation is a new estimate for the number of critical points of the Monge--Amp\`ere operator, which controls the size of the singular set. This estimate uses the \textbf{focal points} of the transport map and the \textbf{cut locus} to partition the domain into regions where the map is well-behaved. In each region, the map is shown to be Lipschitz, and a covering argument bounds the Hausdorff dimension.

The proof also uses the theory of \textbf{convex geometry} and the \textbf{Alexandrov theorem}, which characterizes the second distributional derivatives of convex functions as Radon measures. The Alexandrov theorem provides a weak notion of second derivatives that is well-suited to the study of Monge--Amp\`ere equations. The key estimate is a bound on the number of points where the Alexandrov Hessian has a large determinant, which controls the size of the singular set.

The partial regularity result has several important corollaries. First, it implies that the optimal transport map is continuous almost everywhere, which is a key ingredient in the proof of the uniqueness of optimal transport maps. Second, it provides a bound on the Hausdorff dimension of the singular set, which is useful for the study of the geometry of the space of probability measures. Third, it has applications to the study of the Monge--Amp\`ere equation in the context of differential geometry and meteorology.

\textbf{Higher integrability (2013).} Also with De Philippis, Figalli proved the sharp Sobolev regularity: $u \in W^{2,p}$ if and only if $\det D^2 u \in L^p$ for $1 \le p \le \infty$. This elegant characterization resolved a long-standing question about the relationship between the integrability of the Monge--Amp\`ere determinant and the second derivatives of convex solutions. The proof uses a delicate duality argument and the theory of divergence-measure fields.

The significance of this result lies in its completeness: it provides a full characterization of the integrability properties of solutions to the Monge--Amp\`ere equation. Before this work, only partial results were known---Caffarelli's interior regularity theory gave $C^{1,\alpha}$ regularity under smoothness assumptions on the data, but the sharp $W^{2,p}$ characterization was missing. The result has implications for the study of the Monge--Amp\`ere equation in the context of optimal transport, where the regularity of the transport map is directly related to the integrability of the Monge--Amp\`ere determinant.

The proof technique uses the theory of divergence-measure fields, which provides a framework for studying vector fields whose divergence is a measure. This theory, developed by Chen and Torres, allows the use of distributional methods in settings where classical PDE theory fails. The key estimate is a bound on the total variation of the divergence measure in terms of the $L^p$ norm of the Monge--Amp\`ere determinant, which is then used to deduce the $W^{2,p}$ regularity of the solution.

\textbf{Quantitative isoperimetric inequality (2010).} With Maggi and Pratelli, Figalli proved a sharp quantitative anisotropic isoperimetric inequality with optimal exponent 2: $P_\Phi(E) - P_\Phi(W) \ge C\,\alpha(E)^2$, where $W$ is the Wulff shape and $\alpha(E)$ the asymmetry index. The exponent 2 is optimal, resolving a conjecture of Hall. The proof uses optimal transport and the theory of sets of finite perimeter. The key idea is to use the optimal transport map to compare the given set $E$ with the Wulff shape $W$, and to use the isoperimetric inequality to bound the deficit in terms of the asymmetry.

\textbf{Stability of functional inequalities (2013--2017).} With Carlen, Figalli proved sharp quantitative stability for the logarithmic Sobolev inequality and Gagliardo--Nirenberg inequalities. This yields explicit convergence rates for the critical-mass Keller--Segel equation modeling chemotaxis. The proof uses the Wasserstein gradient flow structure of the Keller--Segel equation and the theory of entropy production. The key insight is that the convergence rate is controlled by the spectral gap of the linearized operator, which can be computed explicitly.

The logarithmic Sobolev inequality states that for a probability measure $\mu$ on $\mathbb{R}^n$,
\[
\int_{\mathbb{R}^n} f^2 \log f^2 \, d\mu \le C \int_{\mathbb{R}^n} |\nabla f|^2 \, d\mu
\]
for all smooth functions $f$ with $\int f^2 \, d\mu = 1$. The equality case is achieved by Gaussian measures, and the stability result quantifies how close a measure must be to a Gaussian if the inequality is nearly an equality. The quantitative stability provides an explicit bound on the Wasserstein distance between the measure and the nearest Gaussian in terms of the deficit in the logarithmic Sobolev inequality.

The application to the Keller--Segel equation is particularly elegant. The Keller--Segel equation describes the aggregation of cells in response to chemical signals, and for critical initial mass, solutions may blow up in finite time. The quantitative stability of the logarithmic Sobolev inequality provides explicit convergence rates for solutions below the critical mass, showing that the cells spread out exponentially fast to the equilibrium distribution. The convergence rate is given by $\lambda_1 t$, where $\lambda_1$ is the first nonzero eigenvalue of the linearized operator, which can be computed explicitly in terms of the Bessel function.

\textbf{Free boundary problems (2015--2019).} Figalli analyzed the optimal partial transport problem, proving regularity of the free boundary using optimal transport techniques and geometric measure theory. The free boundary separates the transported mass from the stationary mass, and its regularity depends on the regularity of the densities and the cost function. The proof uses the \textbf{obstacle problem} and the theory of \textbf{viscosity solutions} to show that the free boundary is a smooth hypersurface.

The connection between optimal transport and the obstacle problem is particularly elegant: the free boundary in the partial transport problem can be characterized as the set where the Kantorovich potential touches a quadratic function, which is precisely the condition that defines the free boundary in the obstacle problem. This observation allows the use of the extensive machinery developed for the obstacle problem, including the theory of thin obstacles and the theory of blow-up limits.

The \textbf{optimal partial transport problem} has applications to several areas of mathematics and science. In image processing, it provides a natural framework for interpolating between two images. In economics, it models the optimal allocation of resources between two competing demands. In biology, it describes the migration of cell populations in response to chemical signals.

\section{Impact of This Discovery}

Figalli's partial regularity theorem is now a standard result in optimal transport. The quantitative stability methods have been extended to Riemannian manifolds, metric measure spaces, and phase-transition models (Allen--Cahn dynamics). The Keller--Segel result resolved a decade-old problem in mathematical biology by providing explicit convergence rates for the critical-mass case.

The construction of moduli spaces of K-stable Fano varieties and the continued development of Monge--Amp\`ere regularity theory on non-convex domains remain active areas influenced by Figalli's work. His lecture notes on optimal transport and Wasserstein gradient flows (2024) provide the modern reference for the field.

Machine learning applications of optimal transport---Wasserstein GANs, diffusion models, and domain adaptation---rely on the theoretical foundations Figalli helped establish. The computational aspects of optimal transport have become increasingly important in data science and machine learning.

The connection between optimal transport and the theory of \textbf{gradient flows} in the Wasserstein space has been developed extensively by Villani, Ambrosio, Gigli, and Savar\'e. Figalli's work on the regularity of optimal transport maps is a key ingredient in the study of gradient flows, which provide a variational framework for studying evolution equations such as the heat equation, the porous medium equation, and the Keller--Segel equation.

The \textbf{Monge--Amp\`ere equation} $\det(D^2 u) = f$ is a fully nonlinear PDE that arises naturally in several contexts. In differential geometry, it appears as the equation for K\"ahler--Einstein metrics: given a K\"ahler form $\omega$ on a complex manifold, a K\"ahler--Einstein metric $\omega_{KE}$ satisfies $\omega_{KE}^n = e^f \omega^n$ for some function $f$, which is equivalent to the complex Monge--Amp\`ere equation. In meteorology, the semi-geostrophic equations reduce to the Monge--Amp\`ere equation via the Monge--Amp\`ere distance. In economics, the equation describes the optimal allocation of resources in a competitive equilibrium.

The regularity theory for the Monge--Amp\`ere equation is one of the central problems in nonlinear PDE theory. Caffarelli's interior regularity theory (1990s) showed that convex solutions are smooth when the data is smooth, but boundary regularity and partial regularity remained open until Figalli's work. The partial regularity result---that the singular set has Hausdorff dimension at most $n-1$---is sharp, as there are examples where the singular set has dimension exactly $n-1$.

The \textbf{quantitative stability} of the isoperimetric inequality has applications to the study of geometric flows. The Allen--Cahn equation $\partial_t u = \Delta u + u - u^3$ describes phase transitions, and its solutions converge to minimal surfaces as a parameter $\epsilon \to 0$. The quantitative stability of the isoperimetric inequality provides bounds on the convergence rate, which are essential for understanding the dynamics of phase transitions. Figalli's work on quantitative stability has been applied to the Allen--Cahn equation, providing sharp convergence rates in various geometric settings.

\section{Modern Views and Improvements}

Figalli continues to work on boundary regularity for optimal transport with non-convex domains and on extending the theory to non-quadratic cost functions. His work on random matrix theory represents a new direction connecting probability and geometry.

The regularity theory for the Monge--Amp\`ere equation continues to be an active area, with recent progress on non-convex domains and rough data. The connections to differential geometry and to the theory of optimal transport maps on manifolds are being developed further.

Recent work of Figalli and collaborators on the structure of the singular set of solutions to the Monge--Amp\`ere equation has revealed new geometric properties. The singular set is conjectured to be countably rectifiable, and partial results in this direction have been obtained using techniques from geometric measure theory.

The \textbf{optimal transportation map} has applications to the study of \textbf{K\"ahler--Einstein metrics} on Fano manifolds, where the Monge--Amp\`ere equation appears as the complex Monge--Amp\`ere equation. The regularity theory developed by Figalli and his collaborators provides tools for studying the existence and uniqueness of K\"ahler--Einstein metrics, which are important in differential geometry and mathematical physics.

The interaction between optimal transport and \textbf{geometric measure theory} continues to be a fruitful area of research. Figalli's work on the regularity of sets of finite perimeter, combined with his work on optimal transport, has led to new results on the structure of minimal surfaces and on the regularity of solutions to geometric evolution equations.

\begin{center}\rule{0.5\linewidth}{0.5pt}\end{center}

\section*{Summary}
\addcontentsline{toc}{section}{Summary}

Figalli proved fundamental regularity results for optimal transport maps and the Monge--Amp\`ere equation, established sharp quantitative stability for geometric and functional inequalities, and applied optimal transport to free boundary problems, the Keller--Segel equation, and meteorology.
